%----------
%	JUSTIFICATION OF THE SOLUTION
%----------	
 
A lo largo de este capítulo, hemos explorado el estado del arte sobre la representación de las mujeres en los medios de comunicación y los sesgos de género presentes en la narrativa periodística. Asimismo, hemos analizado las distintas herramientas de Inteligencia Artificial empleadas para la detección de estos sesgos, comprobando que su integración está reconfigurando los procesos productivos de las redacciones. Sin embargo, esta transformación genera oportunidades a la vez que revela el riesgo de perpetuar desigualdades. Diversos estudios muestran que la representación en la cobertura mediática sigue siendo desigual: más del 75 \% de las fuentes expertas citadas son hombres \cite{unwomen2025gmmp}, un sesgo que los LLM tienden a heredar, asociando nombres femeninos con roles vinculados al cuidado entre tres y seis veces más que los masculinos \cite{kotek2023}.

Ante este panorama, y debido al inmenso volumen de contenido que se publica diariamente, resulta clave contar con metodologías capaces de analizar estos sesgos con rigor y de forma escalable. Sin embargo, como se advierte en \cite{lai2021}, ni los enfoques puramente manuales ni los completamente automatizados pueden abordar por sí solos la complejidad del problema: los primeros son costosos y propensos al error humano, mientras que los segundos pueden generar riesgos de explicabilidad, éticos y legales. Esta limitación ha impulsado el desarrollo de propuestas híbridas que combinan el procesamiento del lenguaje natural, los LLM y la supervisión humana (\textit{human-in-the-loop}), uniendo la potencia del aprendizaje automático con la sensibilidad interpretativa de expertas y expertos en estudios de género y comunicación.

Dentro de este ecosistema de soluciones tecnológicas, y como se ha expuesto en la revisión previa, cabe destacar que, si bien plataformas recientes como \textit{HODIO} suponen un avance analítico sin precedentes para diagnosticar el problema de forma externa y \textit{a posteriori}, su enfoque subraya la necesidad complementaria de herramientas de intervención \textit{a priori}, como la propuesta en este Trabajo de Fin de Máster, que actúen directamente en la propia redacción antes de que el sesgo sea publicado.

Precisamente para dar respuesta a esta necesidad, y en un contexto donde el sector periodístico avanza hacia una integración progresiva de herramientas de IA en los flujos de trabajo editoriales, surge este proyecto, que propone el diseño de un sistema basado en grandes modelos de lenguaje y agentes especializadas que asisten a la analista sin sustituir su criterio. El sistema integra la metodología y el conocimiento experto para identificar y localizar los patrones de sesgo en el texto, y aporta para cada uno las evidencias y una explicación que la profesional revisa y, si procede, corrige, contribuyendo a una comunicación más justa y libre de sesgos.